% =============================================================================
%  CHAPTER 6 — CONCLUSION, REFLECTION & FUTURE WORK
%  SOURCE:   docs/EVALUATION.md §9 (refined hypothesis), §10 (developer study)
%  RUBRIC:   Evaluation & Reflection (20%)
%  BUDGET:   ~4-6 pages
%  TOP BAND: mature self-critical judgement + broader impact analysed with depth.
%
%  >>> DRAFTED HERE (AI first draft, grounded in your docs — verify, then edit):
%        - Summary of Contributions (calibrated to MEASURED results)
%        - Limitations, Future Work (incl. the RQ4 study design from §10)
%  >>> YOURS TO WRITE (\writeme — these are judgement parts, per your working
%      agreement): Reflection, Broader Impact. Scaffolds are provided to react to.
% =============================================================================
\chapter{Conclusion and Reflection}
\label{ch:conclusion}

% --- AI FIRST DRAFT. Edit into your own voice; verify every claim is MEASURED. ---
This project set out to test a single architectural proposition: that an LLM
agent performing refactoring should be \emph{decoupled} from the act of writing
code --- deciding \emph{what} to change while delegating \emph{how} to the IDE's
structured refactoring engine. \Cref{ch:design,ch:impl} realised this as an
IntelliJ plugin that exposes rename, move, change-signature, safe-delete, inline,
extract, and symbol-creation operations through a localhost tool API that any
function-calling model can drive, and \Cref{ch:eval} evaluated it against a
text-editing baseline.

\section{Summary of Contributions}
\label{sec:concl-contributions}
% --- AI FIRST DRAFT. Calibrate each bullet to what you actually MEASURED. ---
The work makes the following contributions, each substantiated earlier in the
report:
\begin{itemize}
  \item \textbf{A structured-refactoring tool API for LLM agents.} An IntelliJ plugin
    (\Cref{ch:design}--\ref{ch:impl}) that exposes IntelliJ's PSI refactoring processors
    through a model-agnostic, function-calling interface (HTTP and MCP), letting an
    external agent \emph{execute} project-wide refactorings that are correct by
    construction rather than emit text. To the best of our knowledge this is the first
    general-purpose interface of its kind, contemporaneous with the JetBrains MCP server.
  \item \textbf{A headless-execution technique for IDE refactorings.} A method for
    driving \texttt{BaseRefactoringProcessor}-based operations --- and editor-bound ones
    such as extract-method --- fully non-interactively (\Cref{ch:impl}), suppressing the
    dialogs and editor dependencies that otherwise make them uncallable by an unattended
    agent.
  \item \textbf{A reproducible benchmark and a controlled comparison.} A compile-anchored
    validation harness (\Cref{sec:eval-validation}) and task suites over three real-world
    Java projects (\Cref{ch:eval}), including a model-held-constant comparison that
    isolates the edit primitive from the driving model's capability.
  \item \textbf{An empirical characterisation of where the action space matters.}
    Evidence (\Cref{sec:eval-rootcause,sec:eval-controlled}) that text-level editing
    fails on reference-graph- and type-dependent refactorings for scripted and
    weaker-model drivers, that a capable agent reaches parity on those refactorings, and
    that the structured engine's durable advantages are its structural queries,
    correctness by construction, and robustness to the driver's capability.
\end{itemize}

\section{Reflection}
\label{sec:concl-reflection}
% --- YOUR SECTION (judgement). The refined hypothesis in EVALUATION.md §9 is
%     genuinely self-critical — that honesty is what earns marks here. ---
The most valuable part of this project, in retrospect, was being proven partly wrong. I
set out expecting the structured engine to be straightforwardly \emph{more correct} than
text editing, and against the scripted and weak-model baselines it is. But the controlled
comparison refined that into something more honest: a capable model edits text correctly
too, even at scale, when it has a compile-and-fix loop, so the structured engine's
durable advantage is not raw correctness but that it is correct \emph{by construction},
supports queries text cannot express, and needs no such loop. Delegating to IntelliJ's
own processors gave that correctness for free on exactly the cases the baselines failed
--- the part of the design that worked best was the part I wrote least of. Starting
again, I would run the cross-model and behavioural (test-based) validation earlier
rather than leaning on compilation and content checks, since those are the experiments
that would most sharpen the claim.

\section{Broader Impact}
\label{sec:concl-impact}
% --- YOUR SECTION (judgement). Required from the 60% band up. ---
A broader implication of this work is that, as capable agents become commonplace, the
differentiator shifts from the models themselves towards the \emph{infrastructure} they
act through --- the structured, verifiable tool layers that let an agent change the
world reliably rather than by best-effort generation. Framed as a safety property, this
is the value of constraining the action space: a small, typed, correct-by-construction
set of operations is safer than an unconstrained ``edit any text'' action, because whole
classes of silently-wrong change simply cannot be expressed
\citep{demystifying2025, sweagent2024}. The benefit is greatest for \emph{less} capable
models, which cannot reason their way around a missing guarantee --- so this kind of
infrastructure is a route to making cheaper, local, or open models usefully reliable,
not only to making frontier models marginally better. The pattern is not specific to
refactoring: the same argument extends to CAD and engineering tooling, to reliable code
generation, and to the widening range of MCP-exposed tools across domains as different
as travel and design.

This shifts what a developer must verify. When the edit itself is performed by a
semantics-preserving engine, attention moves from ``did the change apply correctly?'' to
``was this the right change to ask for?'' --- a genuine reduction in the surface for
silently-wrong edits to reach a codebase. The same property cuts both ways, however. A
guarantee that the mechanics are correct can encourage \emph{over}-trust, with
developers reviewing agent-applied changes less carefully precisely because they compile
and look safe; and it concentrates influence in whoever builds and controls the tool
layer, since every agent that depends on it inherits its blind spots. Infrastructure
that makes automated change safer also makes large automated change easier --- a
responsibility as much as a capability.

Finally, the efficiency dimension carries a sustainability angle: performing a
refactoring in a single structured operation rather than many file reads and per-site
edits reduces the API round-trips, token cost, and therefore the energy spent per task
--- a small but real consideration as agentic workflows scale up.

\section{Limitations}
\label{sec:concl-limitations}
% --- AI FIRST DRAFT. Keep these honest and aligned with sec:eval-threats. ---
The evaluation has several limitations, stated here so that the contributions are
not overread. The benchmark covers three real-world projects but a single language (Java; Kotlin
symbol-creation is implemented but not benchmarked), and the operation set, while
spanning seven families, is not exhaustive. Correctness is established by compilation
and disk-state checks rather than by behavioural (test-suite) equivalence, so a change
that compiles and passes the content checks but alters runtime behaviour would not be
caught. The original expectation --- that the structured approach's correctness
advantage would \emph{grow} on large codebases, where an LLM can no longer hold every
file in context --- was tested directly (on Apache Commons Lang and jackson-databind, up
to 474 files) and was \emph{not} borne out for a capable agent: such an agent reached
parity even at that scale (\Cref{sec:eval-controlled}), so the structured engine's
advantage lies in structural queries and correctness-by-construction rather than in raw
scale. The agent arms were also driven by a single capable model (Claude Opus~4.8);
confirming that the pattern is not specific to one model is left to future work, as is
the human-factors question.

\section{Future Work}
\label{sec:concl-future}
% --- AI FIRST DRAFT from EVALUATION.md §10. The RQ4 study design is yours; this
%     frames it as planned-not-done. ---
Several directions follow directly from the limitations above.

\paragraph{Controlled developer study (RQ4).} The most direct extension is the
human-factors study the evaluation deferred. The planned design compares three
conditions --- manual IntelliJ refactoring, an in-IDE chat assistant performing
text edits, and the structured-refactoring agent --- on a fixed set of maintenance
tasks performed by 10--15 participants on an unfamiliar Java project, measuring
time to completion, the number of corrective undo actions, and post-task cognitive
load (NASA-TLX), with a short interview probing whether participants trusted each
change without manually verifying it. This would test whether the structural
correctness guarantees translate into measurable gains in developer effort and
trust, which the present, purely technical, evaluation cannot establish.

\paragraph{Cross-model studies.} The scale comparison was carried out as part of this
work --- on Apache Commons Lang and jackson-databind, up to 474 files
(\Cref{sec:eval-controlled}) --- so the open question is no longer scale but
\emph{generality across models}. Re-running the capable-agent arm on other frontier
models (for example a GPT- or Gemini-class model, through the same model-agnostic tool
interface) would establish whether the observed pattern --- parity on refactoring
correctness, with the structured advantage concentrated in structural queries and
by-construction safety --- holds beyond Claude, as the architecture intends.

\paragraph{Broader operation and language coverage.} Several higher-level operations
are implemented but not yet benchmarked (pull-up and push-down member,
make-method-static, type migration); benchmarking these, adding the operations still
missing (e.g.\ extract-interface, encapsulate-field), and exercising the existing Kotlin
support would broaden the claim beyond the rename, move, and change-signature core
evaluated here. \paragraph{IDE-agnostic backend.}
Because the agent interacts only through the tool API, the same interface could in
principle be backed by another language service (e.g.\ an LSP server or a
different IDE), which would generalise the architecture beyond IntelliJ.

\paragraph{Stronger evaluation oracles.} Correctness here is established by
compilation and content checks; two complementary oracles would deepen it. First,
structural-quality analysis with DesigniteJava --- the methodology of
\citet{agenticrefactoring2025} --- would quantify whether structured refactorings
preserve or improve design-smell and complexity metrics, making the results
directly comparable to that study (and testing the expectation that
behaviour-preserving operations introduce no new smells). Second, the
differential-fuzzing approach of \citet{difffuzzing2026} would verify behavioural
equivalence beyond compilation --- a behavioural complement to the structural
guarantees the engine already provides by construction.

My own view is that AI-assisted refactoring should become both more widespread and more
\emph{language-aware}. Most agentic coding today happens at the command line and edits
code as text; the natural next step is to wire these agents, by default, into the
structured engine of whatever language they are working in --- the PSI for Java and
Kotlin, an equivalent service elsewhere --- so that whenever the language is known, edits
flow through its abstract syntax tree rather than through free text. I believe that shift
could meaningfully improve how safely and correctly code is written, precisely as more of
it comes to be written with AI.
